%% CV for Aalborg Universitet - Data- og AI-konsulent (Forskningsservice)
%% Compiled with: lualatex

\documentclass[11pt,a4paper,sans]{moderncv}
\moderncvstyle{banking}
\moderncvcolor{blue}

\renewcommand*{\firstnamestyle}[1]{{\fontsize{34}{36}\bfseries\upshape\color{color1}#1}}
\renewcommand*{\lastnamestyle}[1]{{\fontsize{34}{36}\bfseries\upshape\color{color1}#1}}
\renewcommand*{\sectionstyle}[1]{{\sectionfont\color{color1}#1}}

\usepackage[utf8]{inputenc}
\usepackage{hyperref}
\hypersetup{
    colorlinks=true,
    linkcolor=blue,
    filecolor=magenta,
    urlcolor=blue,
    pdftitle={Jacob El-Omar - CV},
    pdfpagemode=FullScreen,
}
\usepackage[scale=0.80]{geometry}
\usepackage{import}

\name{Jacob}{El-Omar}
\address{Aalborg, Denmark}{}{}
\phone[mobile]{+45 53 35 27 82}
\email{elomarjc@gmail.com}
\extrainfo{\href{https://linkedin.com/in/jacob-el-omar}{LinkedIn} \textbullet{} \href{https://github.com/elomarjc}{GitHub}}

\begin{document}

\makecvtitle

% ============================================================
%     PROFILE
% ============================================================

\vspace{6pt}
\small{Civilingeniør (MSc) i Elektroniske Systemer fra Aalborg Universitet med stærk ekspertise inden for dataanalyse, database-infrastruktur og praktisk anvendelse af AI-teknologier. Erfaring med at strukturere spredte datakilder, bygge databaser (SQL, PostgreSQL, Supabase) og afprøve AI-modeller i udviklingsprojekter. Kombinationen af en analytisk profil, dyb indsigt i AAU's organisation og stærke samarbejdsevner gør mig klar til at understøtte Forskningsservices data- og AI-transformation.}

% ============================================================
%     TECHNICAL SKILLS
% ============================================================

\section{Faglige Kompetencer}
\vspace{1pt}
\begin{itemize}

\item \textbf{Data Management \& Databaser}: Datamodellering og -strukturering, databaseintegration (SQL, PostgreSQL, MySQL, Supabase, SQLite), datarensning og ETL-processer.

\item \textbf{AI-Værktøjer \& Maskinlæring}: Praktisk erfaring med implementering af AI-algoritmer i Python, prompt engineering, anvendelse af agentiske AI-værktøjer som Claude Code til automatisering, samt data-visualisering (Matplotlib, Pandas).

\item \textbf{Systemudvikling \& Web}: Backend-udvikling, REST-API-integration, versionsstyring (Git), automatisering via GitHub Actions, samt styring af WordPress platforme.

\item \textbf{Konsulent \& Formidling}: Formidling af komplekse tekniske problemstillinger, tværfagligt samarbejde (PBL-metodik), procesoptimering og udarbejdelse af teknisk dokumentation og beslutningsgrundlag.

\end{itemize}

% ============================================================
%     EDUCATION
% ============================================================

\section{Uddannelse}
\vspace{1pt}

\cventry{Sep 2023--Jun 2025}{MSC i Elektroniske Systemer}{Aalborg Universitet}{Aalborg, Danmark}{}{
Kandidatspeciale: ``Adaptive Noise Cancellation For Electronic Stethoscopes.'' Modellering af komplekse stokastiske signaler, statistisk dataanalyse, datarensning samt algoritmeudvikling og -afprøvning i Python og MATLAB.
}

\vspace{3pt}

\cventry{Sep 2020--Aug 2023}{BSc i Elektronik og IT}{Aalborg Universitet}{Aalborg, Danmark}{}{
Specialisering inden for systemarkitektur, databaseintegrationer, digitalt design og anvendt data-analyse. Projekter omfattede bl.a. reeltidssystemer, robotnavigationsdata og sensorstyring.
}

% ============================================================
%     PROFESSIONAL EXPERIENCE
% ============================================================

\section{Erhvervserfaring}
\vspace{3pt}

\cventry{Sep 2024--Jun 2025}{AI/ML Engineer (Praktikant)}{Ai Health Highway}{Aalborg, Danmark / Remote}{}{
\begin{itemize}
    \item Udviklede, optimerede og testede signalbehandlings- og dataanalysealgoritmer i Python.
    \item Håndterede og strukturerede store medicinske datasæt, rensede måledata og evaluerede algoritmepræstationer i forhold til medicinske standarder.
    \item Udarbejdede teknisk dokumentation og afrapportering af testresultater til brug for tværfaglige, internationale teams.
\end{itemize}}

\vspace{3pt}

\cventry{Jan 2026--Nuværende}{IT- \& Webudvikler (Frivillig)}{Dawah Danmark}{Aalborg, Danmark}{}{
\begin{itemize}
    \item Moderniserede, strukturerede og vedligeholdt databaser og webinfrastruktur (WordPress, MySQL) for at sikre høj oppetid og systemintegritet.
    \item Forbedrede dataoverførsel og systemarkitektur for at optimere interne arbejdsgange og dataopsamling.
\end{itemize}}

\vspace{3pt}

\cventry{Jun 2023--Aug 2024}{Lager- og Logistikmedarbejder}{Lyn Express}{Aalborg, Danmark}{}{
\begin{itemize}
    \item Håndterede registrering, katalogisering og overvågning af tidsfølsomme logistikdata under højt arbejdstempo.
\end{itemize}}

% ============================================================
%     SELECTED PROJECTS
% ============================================================

\section{Udvalgte Projekter}
\vspace{1pt}
\begin{itemize}

\item{\textbf{Grow a Fish (Fritidsprojekt)}: Egenudviklet full-stack mobilspil med vækstdynamikker, multiplayer og in-game butik (Flutter, Dart, SQL, REST API).}

\item{\textbf{AAUSAT6 ADCS Development}: Udvikling af attituderegulatorer til CubeSat-stabilisering ved hjælp af LQR og robust regulering, herunder B-dot-algoritmer (MATLAB, Simulink, LQR, Kalman-filter).}

\item{\textbf{CoM Calibration of 3-DoF Satellite Simulator}: Algoritmeudvikling til estimering af tyngdepunkt (CoM) ved brug af kinematik og Kalman-filtrering (MATLAB, Kalman-filter, Kinematik, OnShape, CNC).}

\item{\textbf{Robot Collision-Avoidance}: Prædiktivt navigationssystem til MiR200-platforme ved brug af UWB og regression (UWB, Kalman-filter, Regression, MiR200).}

\end{itemize}

% ============================================================
%     LANGUAGES
% ============================================================

\section{Sprog}
\vspace{1pt}
\begin{itemize}
\item Dansk (Modersmål), Engelsk (Flydende i skrift og tale).
\end{itemize}

% ============================================================
%     REFERENCES
% ============================================================

\section{Referencer}
\vspace{1pt}
\begin{itemize}
\item Kan oplyses efter anmodning.
\end{itemize}

\end{document}
